\begin{abstract}
Child speech emotion recognition (SER) systems are routinely trained and evaluated on acted corpora, yet deployed in naturalistic settings where vocal expressions are subtle, noisy, and distributionally distinct.
We study this gap through a distribution-driven pipeline built on frozen WavLM features, learnable layer fusion, and a strictly parameter-matched comparison ($N{=}111{,}105$ trainable parameters per pooling head) between prosody-guided temporal importance pooling and pure self-attention pooling.

Across six controlled experiments on acted child speech (C-BESD), spontaneous child speech (FAU Aibo), and adult speech (IEMOCAP), three findings emerge.
First, zero-shot transfer from acted to spontaneous child speech collapses to \textbf{18.70\%} weighted accuracy (WA)---below the four-class chance level---despite \textbf{92.78\%} in-domain acted performance, confirming that distribution shift, not architecture, is the dominant failure mode (Fr\'{e}chet distance FD$=$16.33).
Second, on naturalistic FAU Aibo speech under strong regularization, prosody-guided pooling matches or slightly exceeds self-attention in WA (\textbf{66.36\%} vs.\ 66.18\%), acting as a structural regularizer against environmental noise.
Third, sacrificing only \textbf{1.48} percentage points on acted C-BESD buys interpretability: attention weights correlate strongly with frame-level energy (APC$_{\text{wav}}{=}0.698$), enabling clinically actionable explanations unavailable from black-box self-attention.

We argue that deployable child SER should be evaluated on spontaneous speech and optimized for explainability, not leaderboard accuracy on acted benchmarks alone.
\end{abstract}
